\section{Objetivos, pregunta guía e hipótesis}\label{sec:objetivos}

\subsection{Objetivo general}

Desarrollar un modelo de clasificación que identifique a los clientes de STEG con mayor
probabilidad de fraude detectado, a partir de su historial de facturación entre 2005 y 2019
y de sus atributos, para priorizar las inspecciones de campo.

\subsection{Objetivos específicos}

\begin{enumerate}
  \item Caracterizar los datos de clientes y facturas en cuanto a su calidad, el desbalance
    de clases y la relación entre la longitud del historial y la tasa de fraude.
  \item Construir variables por cliente y comparar modelos de clasificación con una
    validación cruzada agrupada por cliente y un conjunto de prueba sellado, de forma que no
    haya fuga de información.
  \item Evaluar el modelo seleccionado como se usaría en la práctica, midiendo cuántos
    fraudes encuentra una lista de inspección de tamaño fijo frente al azar y frente a una
    línea base de exposición, e interpretar qué variables pesan más en sus predicciones.
\end{enumerate}

\subsection{Pregunta guía}

¿Permite el historial de facturación ordenar a los clientes según su probabilidad de fraude
detectado mejor que el azar y que una referencia basada solo en la exposición a inspección,
una vez se retiran del historial las huellas que deja la propia inspección?

\subsection{Hipótesis}

\begin{description}
  \item[H1.] Una parte del desempeño que se obtiene con el historial completo procede de
    registros posteriores a la detección. Si es así, las variables calculadas sobre la
    segunda mitad del historial rinden claramente más que las de la primera, y esa
    diferencia se reduce al censurar el historial en el primer evento de equipo o de
    contrato.
  \item[H2.] Con el historial censurado, un modelo de boosting sobre variables de
    comportamiento concentra más fraudes en la cabeza de la lista que la línea base de
    exposición, y conserva capacidad de ordenamiento dentro de los clientes que no tuvieron
    ningún evento de equipo o de contrato.
\end{description}

La verificación de H1 quedó registrada como propuesta en la bitácora de decisiones del
proyecto al cerrar el análisis exploratorio, antes de entrenar cualquier modelo.
